


\begin{abstract}
  LLM agents struggle to automate complex enterprise procedures because static, pre-built harness cannot capture proprietary policies, undocumented behaviors, or state-dependent recovery paths that surface only during execution. Procedural memory accrued from agent experience offers a training-free path to closing this gap—yet unlike episodic records or semantic facts, it must consolidate repeated successes and diagnosed failures into reusable actions and recovery strategies. We introduce \method (\(\mathcal{W}\)iki-based \(\mathcal{A}\)gent-\(\mathcal{S}\)ynthesized \(\mathcal{P}\)rocedural \(\mathcal{M}\)emory), a nimble but effective method that distills execution traces and feedback into persistent, interlinked wiki pages, enabling continual agent learning. On TheAgentCompany benchmark, WASP-Mem raises the checkpoint pass rate by \textbf{4.6–6.1}\% over the strongest baseline we compared across three models while requiring \textbf{57–61}\% less storage. We further present a systematic controlled study spanning consolidation method, feedback type, trace ingestion scope, and refinement rounds. Our results reveal that prescriptive knowledge organization such as human defined schemas and explicit deduplication rules hurts performance rather than closing the procedural-knowledge gap. In contrast, human diagnosis converts failed traces into corrective knowledge and drives sustained improvement: a single round of human diagnosis outperforms four rounds of self-judgment by \textbf{7.0–10.9}\%. These findings point to a simple division of labor: let agents organize procedural memory, and invest human effort where it matters most—diagnosis.
\end{abstract}
